\chapter{Streszczenie w języku polskim i angielskim}
\label{cha:streszczenie}

\section{Streszczenie w języku polskim}
\label{sec:streszczeniePL}

Praca opisuje aplikację \emph{Testownik} - program do tworzenia i przeprowadzania testów oraz
zarządzania wynikami uczniów. Celem projektu było danie nauczycielowi narzędzia, które obejmuje cały
cykl pracy z testem: od ułożenia pytań, przez przypisanie testu grupom uczniów, po automatyczne
i ręczne ocenianie odpowiedzi oraz analizę wyników.

Aplikacja została napisana w języku C\# na platformie .NET z wykorzystaniem technologii
\emph{Windows Presentation Foundation} (WPF) do budowy interfejsu graficznego oraz \emph{Entity
Framework Core} jako warstwy dostępu do relacyjnej bazy danych PostgreSQL. System rozróżnia dwie
role użytkowników - nauczyciela i ucznia - udostępniając każdej z nich odrębny zestaw funkcji.
Nauczyciel tworzy testy złożone z pytań jednokrotnego wyboru, wielokrotnego wyboru oraz pytań
otwartych, zarządza grupami uczniów i analizuje wyniki, natomiast uczeń rozwiązuje przypisane mu
testy w ograniczonym czasie i przegląda swoje oceny. W aplikacji zaimplementowano między innymi
limit podejść do testu, automatyczne ocenianie pytań zamkniętych, ręczne ocenianie i edycję oceny,
sterowanie widocznością testów, wykrywanie prób oszustwa, eksport wyników do pliku CSV oraz
haszowanie haseł. Efektem końcowym jest działająca aplikacja okienkowa gotowa do użycia.

\section{Streszczenie w języku angielskim}
\label{sec:streszczenieEN}

\selectlanguage{english}
This document describes the \emph{Testownik} application - an information system for creating,
sharing and conducting tests, as well as for managing students' results. The main goal of the
project was to provide the teacher with a tool that automates the entire workflow of working with a
test: from composing questions, through assigning the test to groups of students, up to automatic
and manual grading of answers and the statistical analysis of results.

The application was written in C\# on the .NET platform using \emph{Windows Presentation
Foundation} (WPF) for the graphical user interface and \emph{Entity Framework Core} as the data
access layer for a relational PostgreSQL database. The system distinguishes two user roles - the
teacher and the student - providing each of them with a separate set of features. The teacher
creates tests consisting of single-choice, multiple-choice and open questions, manages groups of
students and analyses results, while the student solves the assigned tests within a limited time
and reviews their grades. The application implements, among others, a limit on the number of
attempts, automatic grading of closed questions, manual grading and grade editing, control over test
visibility, cheating detection, exporting results to a CSV file and password hashing. The end result
is a working desktop application ready for use.
\selectlanguage{polish}
